\documentclass[11pt]{article}
\usepackage{amsmath,amssymb,geometry}
\geometry{margin=25mm}
\begin{document}
\section*{ROOT V149: Lorentz, helix and fine-structure audit}
\noindent These are conditional mathematical relations, NOT a proof of force unification.\\
Unit convention: $c=1$ only; $\hbar\ne 1$ and $G\ne 1$ remain dimensional.
\begin{align}
\boldsymbol r(t)&=(a\cos\omega t,a\sin\omega t,ut),&a^2\omega^2+u^2&=1,\\
\kappa&=a\omega^2,&\tau&=u\omega,&\kappa^2+\tau^2&=\omega^2,\\
D&=\gamma(1-\beta u),&\omega'&=\omega/D,&u'&=\frac{u-\beta}{1-\beta u},\\
\chi&=\frac{\kappa^2}{\kappa^2+\tau^2}=a^2\omega^2,&\chi'&=\frac{a^2\omega^2}{D^2},\\
\beta&=u\implies\chi'=1,&\tau'&=0,\\
\alpha&=\frac{e_{\mathrm{HL}}^2}{4\pi\hbar},&\alpha_{G,e}&=\frac{Gm_e^2}{\hbar},\\
\boldsymbol E&=A(\cos\varphi,\sin\varphi,0),&
\boldsymbol B&=A(-\sin\varphi,\cos\varphi,0),\quad\varphi=k(t-z),\\
\boldsymbol E\cdot\boldsymbol B&=0,&\boldsymbol B^2-\boldsymbol E^2&=0.
\end{align}
\noindent In SI comparison use $G_{\rm SI}m_{e,\rm SI}^2/(\hbar_{\rm SI}c_{\rm SI})$, solely for conversion.
\end{document}
